\begin{table}[t]
\TBL{\caption{Selected FAOSTAT tea-leaves production indicators. Production is reported on a green-leaf basis. Superscripts: e = estimated by the publishing agency; i = imputed by a receiving agency. Percentages are calculated from the 2024 world total and the 2000 national value; minor differences may result from rounding. Source: FAOSTAT QCL, tea leaves (item 667), with values reproduced from the published 2000--2024 series. \cite{faostatqcl,teahq}.\label{tab:faostat-tea}}}{\begin{tabularx}{\textwidth}{>{\raggedright\arraybackslash}p{27mm} >{\raggedleft\arraybackslash}p{25mm} >{\raggedleft\arraybackslash}p{23mm} X}
\toprule
\textbf{Country} & \textbf{2024 production (t)} & \textbf{World share} & \textbf{Change, 2000--2024}\\
\midrule
China & 16,298,900\textsuperscript{e} & 49.8\% & +434\%\\
India & 6,156,031\textsuperscript{i} & 18.8\% & +71\%\\
Kenya & 2,687,200 & 8.2\% & +162\%\\
Sri Lanka & 1,514,024\textsuperscript{i} & 4.6\% & +14\%\\
Türkiye & 1,414,000 & 4.3\% & +86\%\\
Viet Nam & 1,154,510 & 3.5\% & +284\%\\
\bottomrule
\end{tabularx}}
\end{table}\begin{table}[t]
\TBL{\caption{TeaBioSens dataset structure and published sensory scoring summary. The biochemical workbook contains the underlying measurements; this chapter uses the openly published blend-design and overall-score table for direct visualization and does not infer biochemical values from the score alone. \cite{teabiosensdata,sarmateabiosens}.\label{tab:teabiosens-summary}}}{\begin{tabularx}{\textwidth}{>{\raggedright\arraybackslash}p{37mm} X}
\toprule
\textbf{TeaBioSens feature} & \textbf{Dataset description}\\
\midrule
Dataset observations & 600 experimental data points\\
Blend design & 30 blends; 4 g total per blend\\
Processed tea types & Green, white, oolong, black\\
Sensory panel & 20 semi-trained assessors\\
Input features & 19 biochemical / calculated parameters\\
Overall score & 60.85--82.59 / 100 in the published 30-blend SI table\\
Published clustering & 17 high-grade samples (72--85); 13 low-grade samples (60--68)\\
\bottomrule
\end{tabularx}}
\end{table}\begin{table}[t]
\TBL{\caption{Illustrative museum metadata corpus for tea bowls. Dimensions and rights are preserved in the accompanying CSV; the table foregrounds origin, chronology, medium, and material-cultural relevance. The corpus is purposive rather than representative. \cite{metopenaccess,smithsonianopenaccess}.\label{tab:museum-teabowls}}}{\begin{tabularx}{\textwidth}{p{16mm} p{25mm} p{18mm} p{30mm} X}
\toprule
\textbf{Collection} & \textbf{Origin / date} & \textbf{Record} & \textbf{Material / style} & \textbf{Material-cultural relevance}\\
\midrule
Met & China; 12th--late 13th c. & 17.179.5 & Chayang stoneware; iron-oxide glaze & Links tea practice to Song-period ceramic production.\\
Met & China; 11th--12th c. & 29.100.240 & Jian stoneware; iron-oxide glaze, metal rim & Makes glaze, surface, and vessel form part of tea perception.\\
Met & China; 13th c. & 16.122.5 & Longquan celadon; gold lacquer repairs & Material afterlife records use and later Japanese repair.\\
Met & Korea; first half 17th c. & 36.120.502 & Stoneware; stamped white-slip design & Tea ware explicitly connected to Japanese-market production.\\
Met & Japan; second half 18th c. & 36.120.518 & Korean-style Kyoto ware & Shows cross-regional emulation and material lineage.\\
Met & Japan; 17th c. & 1982.465 & Hagi ware & Demonstrates regional ceramic form within tea culture.\\
Smithsonian & Japan; mid 19th c. & F1896.54 & Kyoto ware; iron glaze & Documents craft, origin, and tea-serving contexts.\\
Smithsonian & Japan; early 19th c. & F1898.462 & Agano ware; clear and iron-brown glazes & Shows local kiln materiality as tea-bowl form.\\
Smithsonian & Japan; 19th c. & F1898.127 & Shigaraki ware; ash glaze & Connects rustic surface qualities to tea aesthetics.\\
\bottomrule
\end{tabularx}}
\end{table}\begin{table}[t]
\TBL{\caption{Comparative Material-Cultural Configurations of Tea in East Asia. Source: Author's synthesis; see especially refs. [9, 10, 12-14, 17, 19, 20].\label{tab:comparative-east-asia}}}{\begin{tabularx}{\textwidth}{>{\raggedright\arraybackslash}p{19mm} >{\raggedright\arraybackslash}X >{\raggedright\arraybackslash}X >{\raggedright\arraybackslash}X}
\toprule
\textbf{Context} & \textbf{Material / practice} & \textbf{Mediating forms} & \textbf{Valuation / identity}\\
\midrule
China (\zh{中国}) & Tea leaf (\zh{茶叶}), processing, water, teaware & Tea treatises, poetry, monastic and literary practice & Knowledge, cultivation, sociability, medicinal and spiritual value\\
Japan (\zh{日本}) & Sencha (\zh{煎茶}), matcha (\zh{抹茶}), bowls, rooms, utensils & Chanoyu / senchado, literati gatherings, visual and material arts & Taste, connoisseurship, lineage, Japaneseness, social distinction\\
Korea (\zh{韩国}) & Tea, vessels, cha-rye (\zh{茶禮}), ritual settings & Royal, Buddhist, literary, domestic and modern national forms & Tradition, ritual continuity, cultural identity\\
Contemporary global market & Specialty leaves (\zh{茶叶}), packaging, brewing equipment & Retail, social media, tasting notes, imagery, lifestyle discourse & Authenticity, craft, provenance, wellness, heritage\\
\bottomrule
\end{tabularx}}
\end{table}\begin{table}[t]
\TBL{\caption{Five-Stage Model of Tea Aestheticization. Source: Author's synthesis.\label{tab:five-stage-model}}}{\begin{tabularx}{\textwidth}{>{\raggedright\arraybackslash}p{9mm} >{\raggedright\arraybackslash}p{28mm} >{\raggedright\arraybackslash}X >{\raggedright\arraybackslash}X}
\toprule
\textbf{Stage} & \textbf{Operation} & \textbf{Material / sensory basis} & \textbf{Cultural effect}\\
\midrule
1 & Differentiation (\zh{差异化}) & Flavor, aroma, color, texture, leaf form, provenance & Differences become perceptible and comparable.\\
2 & Mediation (\zh{中介}) & Water, heat, vessel, timing, gesture, room & Perception is organized through material infrastructures.\\
3 & Stabilization (\zh{稳定化}) & Repeated comparison and description & Preferences become norms, styles, and traditions.\\
4 & Socialization (\zh{社会化}) & Ritual, connoisseurship, education, markets & Taste becomes competence, belonging, status, or identity.\\
5 & Recontextualization (\zh{再语境化}) & New media, markets, heritage narratives, lifestyle forms & Existing tea practices acquire new meanings and uses.\\
\bottomrule
\end{tabularx}}
\end{table}

Such a theory also places literature in a different position than a simple ``reflection'' model would allow. Poems, essays, treatises, diaries, and inscriptions do not merely record tea culture after the fact. They participate in making certain sensory distinctions thinkable and memorable. Conversely, material practices constrain literary language by providing the concrete qualities that writers can perceive, classify, and metaphorize. Text and substance form a reciprocal field.

The same can be said of visual art. A painting of a tea gathering, a manuscript album, a decorated bowl, or a contemporary package does not simply depict a pre-existing cultural meaning. Each object frames attention and selects aspects of tea practice as significant. Visual culture is therefore a technology of valuation.

Ritual occupies a similar position. Ritual does not merely symbolize community or hierarchy; it trains the body in ways of moving, tasting, looking, and waiting. The material specificity of tea gives ritual something to organize around. In turn, ritual increases the cultural legibility of material difference.

Finally, consumption should not be positioned as the endpoint of the process. Consumption is one of the sites where cultural meanings are renewed or revised. Every cup reopens the relation among plant, person, object, and narrative. In this sense, tea has no final cultural form. Its meanings remain dependent on practices of preparation, perception, circulation, and reinterpretation.

\begin{backmatter}
\section{Conclusions}

This chapter has argued that the cultural history of tea is best understood not as the progressive accumulation of symbolic meanings around a stable botanical object, but as an ongoing process of material and sensory mediation. \textit{Camellia sinensis} becomes culturally consequential through transformation: cultivation turns plant life into agricultural practice; processing turns fresh leaves into differentiated materials; brewing converts those materials into sensory events; utensils organize perception; literary and visual forms make distinctions communicable; ritual socializes them; and markets, institutions, and identity projects recode them for new contexts.

The example of Rai San'yo's Sanshi Suimeisho condenses this wider history. In the literati pavilion, sencha is simultaneously a plant-derived substance, a beverage, a sensory experience, a technical practice, an object mediated by utensils, an element of sociability, and a prompt for poetry. Its aesthetic value does not arrive after its material properties have been ignored. It emerges precisely because those properties are attended to, differentiated, and incorporated into a larger cultural practice. The tea gathering is therefore a site where botanical matter becomes literary and social form without losing its material specificity.
The broader East Asian comparison reinforces this point. Chinese tea poetry and the \textit{Classic of Tea} reveal how sensory and technical distinctions entered literary and knowledge systems; Japanese tea history demonstrates how utensils, spaces, ritual movements, and connoisseurship constituted a complex aesthetic environment; Korean tea history shows how practices could be reconfigured through modern projects of tradition and national identity. Across these cases, tea's meanings are inseparable from the material techniques through which it is cultivated, processed, prepared, handled, tasted, and displayed.

The concept of aestheticization developed here consequently differs from a purely symbolic account of culture. To aestheticize tea is to make material difference socially perceptible and culturally consequential. It is to transform qualities of matter into objects of attention, objects of attention into standards of taste, and standards of taste into practices of distinction, belonging, and identity. The process is recursive: cultural narratives shape how tea is perceived, while material encounters can confirm, alter, or resist those narratives.

This material-cultural perspective also helps explain why tea remains unusually durable as a cultural form. Unlike an image or abstract symbol, tea must be repeatedly enacted. The leaf continues to demand preparation. Water must meet it; vessels must contain it; bodies must drink it. Every new context therefore reopens possibilities for reinterpretation. Tea can remain traditional while changing, intimate while commercial, ordinary while prestigious, and local while globally circulated.

To follow the leaf across these transformations is ultimately to see that aesthetic culture is not produced by minds alone. In this respect, Bennett's account of vibrant materiality helps sharpen the chapter's claim that cultural processes emerge through configurations in which nonhuman material forces participate without being treated as human-like intention \cite{bennett}. It is generated through encounters among plants, substances, tools, bodies, texts, images, institutions, and environments. The cultural history of \textit{Camellia sinensis} is therefore also a history of how humans learn to perceive matter---and how matter, through its affordances and resistances, helps shape the cultural forms through which humans perceive themselves. The chapter's contextualization of Rai San'yo and Sanshi Suimeisho also draws on publicly available Yale archival and lecture materials \cite{yale}.

\section*{Acknowledgments}
The author thanks the Council on East Asian Studies at Yale University for making publicly available information about the lecture and archival materials that informed the discussion of Rai San'yo and sencha.

\section*{Conflict of interest}
The author declares no conflict of interest.

\section*{Author details}
\begin{authordetails}
\author{Vivien Jiaqian Zhu}
\address[1]{Stanford University, Stanford, California, United States}
\address{*Address all correspondence to: jvzhu@stanford.edu}
\IntechOpentext{\textcopyright\ \the\year{} The Author(s). License and copyright statement to be completed during IntechOpen production.}
\end{authordetails}

\section*{References}

\begin{thebibliography}{99}
\bibitem{shimamura} Shimamura Y. \textit{Rai Sanyo to sencha: Kinsei koki no bunjin no shumi to sono seishinsei ni kansuru shiron} [\zh{頼山陽と煎茶：近世後期の文人の趣味とその精神性に関する試論}]. Tokyo: Kasama Shoin; 2022. ISBN 9784305709585.
\bibitem{brown} Brown TG. \textit{Laws of the Land: Fengshui and the State in Qing Dynasty China}. Princeton, NJ: Princeton University Press; 2024. DOI: 10.1515/9780691246727.\n\bibitem{zhucoetzee} Zhu VJ. \textit{The Visual Ethnographer's Data Other: Secrets Unveiled for a Sociological J. M. Coetzee}. Eliva Press; 2026.
\bibitem{zhusinojapanese} Zhu VJ. \textit{Sino-Japanese Literature in Perspective: A Short Communication to the World Literature} (\zh{中日文学透视：世界文学的短程通信}). Edited by Irina Lungu. Chisinau: Eliva Press; 2026.
\bibitem{appadurai} Appadurai A, editor. \textit{The Social Life of Things: Commodities in Cultural Perspective}. Cambridge: Cambridge University Press; 1986. DOI: 10.1017/CBO9780511819582.
\bibitem{kopytoff} Kopytoff I. The cultural biography of things: commoditization as process. In: Appadurai A, editor. \textit{The Social Life of Things: Commodities in Cultural Perspective}. Cambridge: Cambridge University Press; 1986. p. 64--92. DOI: 10.1017/CBO9780511819582.004.
\bibitem{ingoldmaterials} Ingold T. Materials against materiality. \textit{Archaeological Dialogues}. 2007;14(1):1--16. DOI: 10.1017/S1380203807002127.
\bibitem{ingoldmaking} Ingold T. \textit{Making: Anthropology, Archaeology, Art and Architecture}. London: Routledge; 2013.
\bibitem{clunas} Clunas C. \textit{Superfluous Things: Material Culture and Social Status in Early Modern China}. Honolulu: University of Hawai'i Press; 1991. DOI: 10.1515/9780824841300.
\bibitem{howes} Howes D, editor. \textit{The Varieties of Sensory Experience: A Sourcebook in the Anthropology of the Senses}. Toronto: University of Toronto Press; 1991.
\bibitem{classen} Classen C. \textit{Worlds of Sense: Exploring the Senses in History and Across Cultures}. London: Routledge; 1993.
\bibitem{darnton} Darnton R. What is the history of books? \textit{Daedalus}. 1982;111(3):65--83.\bibitem{chartier} Chartier R. \textit{The Order of Books: Readers, Authors and Libraries in Europe between the Fourteenth and Eighteenth Centuries}. Translated by Lydia G. Cochrane. Stanford, CA: Stanford University Press; 1994.
\bibitem{turner} Turner VW. \textit{The Ritual Process: Structure and Anti-Structure}. Chicago: Aldine Publishing Company; 1969.
\bibitem{bell} Bell C. \textit{Ritual Theory, Ritual Practice}. New York: Oxford University Press; 1992.
\bibitem{jauss} Jauss HR. \textit{Toward an Aesthetic of Reception}. Translated by Timothy Bahti. Minneapolis: University of Minnesota Press; 1982.
\bibitem{smithheritage} Smith L. \textit{Uses of Heritage}. London; New York: Routledge; 2006.
\bibitem{bennett} Bennett J. \textit{Vibrant Matter: A Political Ecology of Things}. Durham, NC: Duke University Press; 2010.
\bibitem{nishimura} Nishimura M, Otsuki Y, Shinohara H, Okamoto H, Miyake M, Miyajima J, Kumano H, Hino Y, Sato M. Tea in the Historical Context of East Asia: Cultural Interactions across Borders. Translated by Jenine Heaton. In: \textit{Cultural Reproduction on its Interface: From the Perspectives of Text, Diplomacy, Otherness, and Tea in East Asia}. The International Academic Forum for the Next Generation Series, vol. 1. Osaka: Institute for Cultural Interaction Studies, Kansai University; 2010. p. 195--226.
\bibitem{benn} Benn JA. \textit{Tea in China: A Religious and Cultural History}. Honolulu: University of Hawai'i Press; 2015. DOI: 10.21313/hawaii/\allowbreak9780824839635.001.0001.
\bibitem{faostatqcl} Food and Agriculture Organization of the United Nations. Crop production, yield, harvested area (Global, National -- Annual) -- FAOSTAT. FAO Data Catalog; revised 18 March 2026. Coverage 1961--2024. Licence: CC BY 4.0. Available at: \url{https://data.fao.org/catalog/iso/d24a448b-3b62-4c09-8c1d-4a39bb599876}.
\bibitem{teahq} TeaHQ. The tea economy, in figures: FAOSTAT tea-leaves series, item 667, 1961--2024. Available at: \url{https://teahq.app/economy}.
\bibitem{bennpoetry} Benn JA. Tea poetry in Tang China (\zh{唐代茶诗}). In: \textit{Tea in China: A Religious and Cultural History}. Honolulu: University of Hawai'i Press; 2015. p. 72--95. DOI: 10.21313/hawaii/\allowbreak9780824839635.003.0004.
\bibitem{teabiosensdata} Sarma O, Dashora K. Tea Biochemical \& Sensory Dataset (TeaBioSens). Zenodo; 2025. DOI: 10.5281/zenodo.17921537.
\bibitem{sarmateabiosens} Sarma O, Dashora K. Machine learning-based prediction of sensory quality in tea blends using a semi-trained panel assessment. \textit{Sustainable Food Technology}. 2026;4:2805--2818. DOI: 10.1039/D5FB00580A.
\bibitem{sang2011} Sang S, Lambert JD, Ho CT, Yang CS. The chemistry and biotransformation of tea constituents. \textit{Pharmacological Research}. 2011;64(2):87--99. DOI: 10.1016/j.phrs.2011.02.007.
\bibitem{lee2019} Lee MK, Kim HW, Lee SH, Kim YJ, Asamenew G, Choi J, Lee JW, Jung HA, Yoo SM, Kim JB. Characterization of catechins, theaflavins, and flavonols by leaf processing step in green and black teas (\textit{Camellia sinensis}) using UPLC-DAD-QToF/MS. \textit{European Food Research and Technology}. 2019;245(5):997--1010. DOI: 10.1007/s00217-018-3201-6.
\bibitem{ho2015} Ho CT, Zheng X, Li S. Tea aroma formation. \textit{Food Science and Human Wellness}. 2015;4(1):9--27. DOI: 10.1016/j.fshw.2015.04.001.
\bibitem{zhang2020} Zhang L, Cao QQ, Granato D, Xu YQ, Ho CT. Association between chemistry and taste of tea: A review. \textit{Trends in Food Science \& Technology}. 2020;101:139--149. DOI: 10.1016/j.tifs.2020.05.015.
\bibitem{juneja1999} Juneja LR, Chu DC, Okubo T, Nagato Y, Yokogoshi H. L-theanine---a unique amino acid of green tea and its relaxation effect in humans. \textit{Trends in Food Science \& Technology}. 1999;10(6--7):199--204. DOI: 10.1016/S0924-2244(99)00044-8.
\bibitem{mair} Mair VH, Hoh E. \textit{The True History of Tea}. London; New York: Thames \& Hudson; 2009.
\bibitem{huang} Huang X. Smellscapes of Nanjing Road: Cognitive and affective mapping. In: Wu S, Huang X, editors. \textit{Sensing China: Modern Transformations of Sensory Culture}. London: Routledge; 2022. p. 71--98. DOI: 10.4324/9781003176220-6.
\bibitem{varley} Varley HP, Kumakura I, editors. \textit{Tea in Japan: Essays on the History of Chanoyu}. Honolulu: University of Hawai'i Press; 1989. DOI: 10.1515/9780824844394.
\bibitem{pitelka} Pitelka M, editor. \textit{Japanese Tea Culture: Art, History and Practice}. London: Routledge; 2013.
\bibitem{kumakura} Kumakura I. \textit{Japanese Tea Culture: The Heart and Form of Chanoyu}. Translated by Martha J. McClintock. Tokyo: Japan Publishing Industry Foundation for Culture; 2023. DOI: 10.2307/jj.2840648.
\bibitem{wuhung} Wu H. \textit{Transience: Chinese Experimental Art at the End of the Twentieth Century}. Revised ed. Chicago: Smart Museum of Art, University of Chicago; 2005.
\bibitem{leipenghan} Lei Y, Peng M, Han Y. Mistaken presuppositions and the translation of cultural texts: a study of two English translations of \textit{The Classic of Tea}. \textit{Perspectives: Studies in Translation Theory and Practice}. 2026;34(4):833--848. DOI: 10.1080/0907676X.2025.2599854.
\bibitem{corbett} Corbett R. Arts of the Tea Ceremony. In: \textit{Oxford Bibliographies in Art History}. Oxford: Oxford University Press; 2025. DOI: 10.1093/obo/9780199920105-0187.
\bibitem{koreanteaclassics} Yi Mok H, Cho-ui. \textit{Korean Tea Classics by Hanjae Yi Mok and the Venerable Cho-ui}. Translated by Hong Kyeong-hee, Brother Anthony, and Steven D. Owyoung. Seoul: Seoul Selection; 2011. ISBN 9788991913660.
\bibitem{xu} Xu H. Translating in and out of East Asian cultures: focus on Chinese and Korean. \textit{Perspectives: Studies in Translation Theory and Practice}. 2024;32(4):569--584. DOI: 10.1080/0907676X.2024.2368038.
\bibitem{carriger} Carriger ML. Tea cultures of Japan. In: Goldstein D, editor-in-chief. \textit{Oxford Research Encyclopedia of Food Studies}. Oxford: Oxford University Press; 2026. DOI: 10.1093/9780197852538.003.0250.
\bibitem{bourdieu} Bourdieu P. \textit{Distinction: A Social Critique of the Judgement of Taste}. Translated by Richard Nice. Cambridge, MA: Harvard University Press; 1984.
\bibitem{hellyer} Hellyer R. \textit{Green with Milk and Sugar: When Japan Filled America's Tea Cups}. New York: Columbia University Press; 2021. DOI: 10.7312/hell19910.
\bibitem{metopenaccess} The Metropolitan Museum of Art. Open Access collection and object records, including tea bowls. New York: The Metropolitan Museum of Art. Available at: \url{https://www.metmuseum.org/art/collection/}.
\bibitem{smithsonianopenaccess} Smithsonian Institution. Open Access collections and object records, National Museum of Asian Art / Freer Gallery of Art. Washington, DC: Smithsonian Institution. Available at: \url{https://www.si.edu/openaccess}.
\bibitem{oshikiri} Oshikiri T. The shogun's tea jar: ritual, material culture, and political authority in early modern Japan. \textit{The Historical Journal}. 2016;59(4):927--945. DOI: 10.1017/S0018246X1600008X.
\bibitem{sen} Sen Soshitsu XV. \textit{The Japanese Way of Tea: From Its Origins in China to Sen Rikyu}. Translated by V. Dixon Morris. Honolulu: University of Hawai'i Press; 1998. DOI: 10.1515/9780824864804.
\bibitem{leekwon} Lee HC, Kwon SH. The invention and promotion of cha-rye (\zh{茶禮}) in Korea. \textit{Asian Journal of Social Science}. 2022;50(3):171--177. DOI: 10.1016/j.ajss.2022.06.004.
\bibitem{decertau} de Certeau M. \textit{The Practice of Everyday Life}. Volume 1. Translated by Steven Rendall. Berkeley: University of California Press; 1984.
\bibitem{yale} Yale Council on East Asian Studies. Steeped in Art and Nature: On Rai San'yo's ``Sanshi Suimeisho''. Yale University; 2026.
\end{thebibliography}
\end{backmatter}

\end{document}